\chapter*{Введение}
\addcontentsline{toc}{chapter}{Введение}

\setbeforesecskip{0.8\baselineskip}
\setaftersecskip{0.35\baselineskip}
\setbeforesubsecskip{0.55\baselineskip}
\setaftersubsecskip{0.25\baselineskip}
\setbeforesubsubsecskip{0.4\baselineskip}
\setaftersubsubsecskip{0.2\baselineskip}
\renewcommand{\chaptername}{Глава}
\renewcommand{\appendixname}{Приложение}
\renewcommand{\listfigurename}{Список рисунков}
\renewcommand{\listtablename}{Список таблиц}
\renewcommand{\bibname}{Список литературы}

\paragraph*{Актуальность темы исследования.}
Российская система территориального планирования связывает доступность объектов обслуживания с показателями национальных программ и проектов. Национальная цель по достижению комфортной и безопасной среды для жизни поставлена в указе Президента РФ №309 от 7 мая 2024 года \cite{rfDecree309_2024}. Национальные направления развития выражаются в федеральных проектах, в частности, в проекте по формированию комфортной городской среды \cite{federalProjectFkgsPassport}. Прикладную часть программы федерального уровня реализуют региональные программы, которые задают уже тот уровень, на котором локальный исполнительный орган власти непосредственно отвечает за достижение целевых показателей эффективности \cite{spbRegionalFkgs}.

Несоответствие заключается в том, что оценки доступности и обеспеченности по нормативам отвечают на вопрос, где и что должно быть размещено, а ПКРТИ отвечает на вопрос, как развивать дорожно-транспортную инфраструктуру. В случае, если эти решения готовятся последовательно, возникает риск создания неэффективного плана застройки территории объектами обслуживания населения, поскольку согласование транспортного плана будет происходить на последующих этапах.

Задача обеспечения доступности объектов предложения не нова. Изначально, города строились с ориентацией на пешеходное движение, затем, с появлением всё новых видов мобильности, городская топология подстраивалась под них. Так, с увеличением расстояний и спроса на личный транспорт, городское пространство стало менее человеко-ориентированным, и, как следствие, менее доступным.


\paragraph*{Степень разработанности темы.}

Классическая линия моделей оптимального размещения объектов формализовала обслуживание как выбор узлов или площадок при заданном спросе и заданной сети расстояний. Задача размещения складов задала базовую постановку выбора складов с учетом затрат обслуживания спроса \cite{kuehn1963warehouse}. P-медиана и p-центр закрепили две ключевые логики: минимизацию суммарного расстояния и ограничение худшего расстояния до объекта \cite{hakimi1964medians}. Для общественных сервисов важным развитием стали покрывающие постановки. Покрытие множества минимизирует число объектов, достаточных для покрытия спроса в пределах нормативного расстояния \cite{toregas1971emergency}. Максимальное покрытие максимизирует покрытый спрос при ограниченном числе объектов \cite{church1974mclp}. В этих и других моделях оптимального размещения транспортная сеть обычно выражена фиксированной матрицей расстояний или времен, а не как самостоятельный объект решения.

Класс задач совместного размещения объектов и проектирования сети включает в себя одновременный выбор мест размещения объектов обслуживания и изменения транспортных связей  \cite{daskin1993towardIntegrated} \cite{melkote1996integrated}. Решение задач покрытия в такой постановке показывает, что удовлетворение спроса зависит от изменений в сети \cite{melkote1998maximumCovering}. Решение этой задачи было также рассмотрено с учетом ограничения на вместимость входящих потоков и различной стоимостью добавления ребер в сети \cite{melkote2000capacitated}. Таким образом, совместное рассмотрение задачи оптимального размещения объектов и моделирования изменений сети был выделено в отдельную категорию Facility-location Network Design problem (FLND)\cite{melkote2001integrated}. Необходимо отметить, что при применении оптимизации размещения объектов к задаче покрытия спроса на доступ к социальной инфраструктуре в большинстве случаев необходимо рассматривать расширенную постановку задачи Facility location с учетом вместимости объектов обслуживания и размера спроса в заданной точке пространства.

\paragraph*{Цель исследования.}
Увеличение эффективности размещения объектов обслуживания населения в рамках процесса территориального планирования

\paragraph*{Задачи исследования.}
Исследование системы управления территориальным планированием в задаче размещения объектов социального обслуживания населения;
Разработка метода оценки транспортной доступности на основе критерия сезонной устойчивости;
Разработка метода оптимизации сети в задаче оптимального размещения;
Применение методов на примере агломераций Арктической зоны РФ.

\paragraph*{Объект исследования.}
Подсистема территориального планирования в составе системы управления развитием территорий

\paragraph*{Предмет исследования.}
Методы и алгоритмы поддержки принятия решений по обеспечению транспортной доступности объектов обслуживания

\paragraph*{Методы исследования.}
В исследовании использованы методы системного анализа, математического моделирования, сетевого анализа, пространственного анализа геоданных, оптимизации размещения объектов обслуживания, оптимизации транспортной сети, эволюционной оптимизации, вычислительного эксперимента и программной реализации алгоритмов поддержки принятия решений.

\paragraph*{Научная новизна.}
Научная новизна исследования состоит в разработке метода оценки транспортной доступности объектов обслуживания с учетом стохастической динамики многоуровневой транспортной сети, позволяющего учитывать изменение состояния сети под воздействием внешних факторов и оценивать сезонную устойчивость сервисной доступности.

Разработан метод интеллектуальной поддержки принятия решений оптимального размещения объектов обслуживания, отличающийся учетом изменяемости транспортной сети и применением оптимизации сетевых параметров как альтернативы или дополнения к размещению новых объектов обслуживания.

Предложена постановка, в которой транспортная сеть рассматривается не только как фиксированный входной параметр задачи размещения, но и как управляемый компонент, изменение которого влияет на требуемое количество объектов обслуживания и уровень обеспеченности населения.

Разработанные методы апробированы в вычислительных экспериментах на реальных пространственных данных для территорий с выраженной зависимостью доступности от внешних факторов и сезонных изменений транспортной сети.


\paragraph*{Теоретическая значимость.}

Теоретическая значимость состоит в развитии теории оценки транспортной доступности за счет учета влияния внешних факторов, качества составляющих сети, учета её конфигурации и учета эффекта взаимосвязи между слоями сети. За счет предложенных расширений получается учесть как изменение состояний самой сети, так и оценить эффект от направленного изменения сети.

\paragraph*{Практическая значимость.}
Практическая значимость работы состоит в программной реализации всех рассматриваемых методов, апробации методов на множестве территорий, что доказывает воспроизводимость вычислений и возможность применения разработанных методов для различных практических приложений.

\paragraph*{Положения, выносимые на защиту.}
\begin{enumerate}
    \item Метод оценки транспортной доступности на основе критерия сезонной устойчивости, отличающийся применением методов сетевого анализа, обеспечивающий учет стохастической динамики многоуровневой сети. Положение соответствует пункту 2 паспорта специальности.

    \item Метод интеллектуальной поддержки принятия решений оптимального размещения объектов обслуживания, отличающийся применением оптимизации сети, обеспечивающий уменьшение необходимого количества объектов обслуживания. Положение соответствует пункту 9 паспорта специальности.
\end{enumerate}

\paragraph*{Достоверность результатов.}
Степень достоверности научных достижений
подтверждается корректным использованием методов, обоснованием постановки задач, экспериментальными исследованиями на реальных массивах данных, покрывающими разработанные технологии и алгоритмы.

\paragraph*{Апробация результатов.}
Результаты исследований обсуждались на конференциях:​

\begin{enumerate}
	\item 40-я ежегодная конференция AAAI по искусственному интеллекту (2026) --- 2-й семинар по ИИ для городского планирования, 2026, рейтинг CORE: A*, стендовый и устный доклад;
	\item Научная конференция AMS (2026) в области городского планирования, специальная секция, устный доклад;
	\item Международная школа и конференция по науке о сетях (NetSci 2026), основная программа, стендовый доклад;
	\item Международная конференция по вычислительной науке и ее приложениям (ICCSA 2023/2024), устный доклад.
\end{enumerate}

\paragraph*{Публикации автора по теме диссертации.}
Основные результаты по теме диссертации представлены в публикациях автора:


\theScopusPapers~опубликована в журнале, индексируемом в базе Scopus.
%One certificate of state registration of a computer program has also been obtained.



Публикации в международных журналах, индексируемых в базе Scopus:
\begin{refsection}[/Users/gk/Code/super-duper-disser/itmo-phd-thesis-template-en/biblio/own.bib]
\nocite{*}
\printbibliography[
    keyword=scopus,
    %title={В международных изданиях, индексируемых в базе данных Scopus},
    %heading=subbibliography,
    heading=none,
    resetnumbers=true
]
\end{refsection}


Список всех публикаций автора по теме диссертации:
\begin{refsection}[/Users/gk/Code/super-duper-disser/itmo-phd-thesis-template-en/biblio/own.bib]
\nocite{*}
\printbibliography[
    keyword=own,
    %title={Список всех публикаций автора по теме диссертации},
    %heading=subbibliography,
    heading=none,
    resetnumbers=true
]
\end{refsection}


\paragraph*{Личный вклад автора.}
В работах, выполненных в соавторстве, вклад автора заключается в
построении моделей, разработке и реализации методов и алгоритмов, разработке, проведении вычислительных экспериментов и подготовке обзоров литературы.

В совместной статье с Губиной А. автор выполнил обзор методов моделирования многоуровневых сетей, обеспечил программную реализацию и разработал алгоритмы расчетов доступности при сезонном изменении в сети.

Во всех работах, выполненных в соавторстве с Тихевич В., автор выполнил постановку задач, методологический обзор и разработку предложенных алгоритмов.

В работах, выполненных совместно с Морозовым А.С. автор выполнил методологическую постановку экспериментов, программную реализацию, а также анализ полученных результатов.

\paragraph*{Структура и объем диссертации.}
Во введении сформулированы цель и задачи
исследования, обоснованы актуальность и научная новизна работы. Также, во введении перечислены основные положения, выносимые на защиту диссертационной работы; представлены практическая и теоретическая значимость.

В первой главе предоставлен литературный обзор по теме диссертации. Показаны состояние предметной области, а также недостатки существующих подходов.

Во второй главе описан метод оценки транспортной доступности 
на основе критерия сезонной устойчивости, описаны результаты применения методов сетевого анализа, показаны результаты учета динамики сезонного изменения сети.

В третьей главе описан метод интеллектуальной поддержки принятия решений размещения объектов обслуживания; описаны основные отличия оптимизации размещения с учетом изменчивости сети.

В четвертой главе описаны экспериментальные исследования эффективности предложенных методов для прикладных задач управления в системе управления развитием территории.


\paragraph*{Соответствие специальности 2.3.4 (в порядке степени соответствия).}

\begin{itemize}
	\item[2.] Разработка математических моделей и критериев эффективности, качества и надёжности организационных систем;
	\item[9.] Разработка методов и алгоритмов интеллектуальной поддержки принятия управленческих решений в организационных системах.
\end{itemize}
